% Statement and cited-input packet. Target proof is separately embargoed.
% Frozen II source SHA256: 74497bfe327fae25c9cdcb17334b3bf42a9793a9df7c0a4404d97568f2f760f0
% Target owner: cdp:st:D.2.2; tag: 01G9

\begin{proposition}[Proved existence and exact Euler number]\hypertarget{tag-01G9}{}\label{cdp:st:D.2.2}


The class of compact complex threefolds with $b_1=b_2=0$,
a proper connected-fibre map with smooth two-torus general fibre,
and trivial rational monodromy is nonempty. Every member satisfies
\[
                   c_3(Y)=0,\qquad b_3(Y)=2.
\]
Every individual reduced fibre support has Euler number zero and
the geometric description of \SeriesRef{cdp:st:C.2.G.1}. Arbitrary singularities and
full-fibre multiplicities are allowed. There is an example with
integral $H_1(Y,\mathbb Z)=H_2(Y,\mathbb Z)=0$ and
algebraic dimension one.
\end{proposition}

% Cited input cdp:st:D.2.3; statement or labelled formula only.
\begin{proposition}[Proved logarithmic period obstruction]\hypertarget{tag-01GA}{}\label{cdp:st:D.2.3}



Let $h:U\to\Delta$ be a proper holomorphic map from a smooth
complex threefold, with connected fibres, smooth two-torus fibres
over $\Delta^*$, and central fibre $F=\sum_i m_iD_i$ with
normal-crossings reduced support $D=\bigcup_iD_i$.
If every normalization component of $D$ has geometric genus zero,
the monodromy on $H^2(U_t;\C)$ is nontrivial. In particular the
monodromy on $H^1(U_t;\Q)$ is nontrivial. No K\"ahler assumption
on $U$ or $D$ is made, and the positive integers $m_i$ are arbitrary.
\end{proposition}

% Cited input cdp:st:D.2.5; statement or labelled formula only.
\begin{proposition}[Proved unrestricted sign under trivial monodromy]\hypertarget{tag-01GC}{}\label{cdp:st:D.2.5}



Let $f:Y\to\PP^1$ be a proper surjective holomorphic map with
connected fibres from a compact connected complex threefold with
$b_1(Y)=b_2(Y)=0$. Suppose its smooth general fibre is a complex
two-torus and its monodromy representation on $H_1(F,\Q)$ is
trivial. Then
\[
 c_3(Y)\le0,\qquad b_3(Y)\ge2.
 \tag{D.S.1}\hypertarget{tag-019O}{}\label{cdp:eq:D.S.1}
\]
The fibres may have arbitrary singularities and arbitrary
multiplicities. In particular this proves the sign assertion of
D.2 for a holomorphic algebraic reduction; no existence assertion
is required for the conclusion.
\end{proposition}

% Cited input cdp:st:D.2.7; statement or labelled formula only.
\begin{proposition}[Proved existence]\hypertarget{tag-01GE}{}\label{cdp:st:D.2.7}



There exists a smooth compact connected complex threefold $Y$
with a holomorphic map $f:Y\to\PP^1$ with connected fibres such
that
\[
 H_1(Y;\Z)=H_2(Y;\Z)=0,\qquad
 a(Y)=1,\qquad c_3(Y)=0,\qquad b_3(Y)=2.
\tag{D.21}\hypertarget{tag-0198}{}\label{cdp:eq:D.21}
\]
Every smooth fibre is the fixed abelian surface
$A=E\times E$, where $E=\C/(\Z+\mathrm i\Z)$, and the marked
integral monodromy on $H_1(A;\Z)$ is the identity. There are
exactly three singular fibres. Each is reduced with simple
normal crossings, and consists of a primary Kodaira surface
and a primary Hopf surface meeting along one smooth elliptic
curve. Consequently the trivial-monodromy class in D.2 is
nonempty, and equality in the unrestricted sign theorem
\ref{cdp:st:D.2.5} is attained.
\end{proposition}

